% GENERATED FILE. Edit canonical lessons, then run npm run generate:books.
% canonical-source-hash: fnv1a64:a7aa803120a7580c
% canonical-lessons: RU-C340-soglashatsya, RU-C340-soglasie, RU-C340-nesoglasie, RU-C340-vozrazhat, RU-C340-vozrazhenie

\chapter{Agreeing and objecting}
\label{ch:ru-agreeing-and-objecting}
\hlchaptermodality{340}

\begin{chapteropening}
\textbf{By the end of this chapter:} \emph{I can agree and give my consent, say that I disagree, and object to an idea.}

Complete the last lesson of chapter 340: I can agree and give my consent, say that I disagree, and object to an idea.
\end{chapteropening}

\section[\texorpdfstring{\ru{соглашаться} (soglashátsya)}{соглашаться (soglashátsya)}]{\ru{соглашаться} (soglashátsya) — to agree}

\label{lesson:RU-C340-soglashatsya}



\begin{quote}
Before the new one: say the Russian for confidence, then the Russian for to doubt.
\end{quote}

\subsection*{You'll want to know: \ru{соглашаться}}
\textbf{\ru{соглашаться}} (\emph{soglashátsya}) — \textquotedblleft{}to agree\textquotedblright{}.

\begin{quote}
\ru{Я} \ru{не} \ru{всегда} \ru{соглашаюсь} \ru{с} \ru{отцом}.
\end{quote}

\begin{quote}
\emph{Ya ne vsegdá soglasháyus s ottsóm.}
\end{quote}

\begin{quote}
I do not always agree with my father.
\end{quote}

Imperfective; perfective \emph{\ru{согласиться}}. \emph{\ru{Соглашаться} \ru{с}} + instrumental (with a person or a view), \emph{\ru{на}} + accusative (to accept an offer).

\begin{cousinweb}
\textbf{\ru{соглашаться}} is built on \emph{\ru{глас}} \textquotedblleft{}voice\textquotedblright{}, the Church Slavonic form of native \emph{\ru{голос}}; both come from Proto-Slavic \emph{*golsŭ} \textquotedblleft{}voice\textquotedblright{} (EtymDB). The \emph{-\ru{ла}-} is the Church Slavonic shape where Russian has \emph{-\ru{оло}-}: to agree is to speak with one voice.
\end{cousinweb}

\begin{grammarlens}[title={What you've built}]
One more word for telling, explaining and describing at B1.
\end{grammarlens}

\subsection*{Your turn}
\begin{itemize}
\raggedright
  \item \emph{Say it:} \emph{soglashátsya}
  \item \emph{Say it:} \emph{soglashátsya}, once more
  \item \emph{Say it:} a sentence of your own with \emph{soglashátsya}
  \item \emph{Recall:} say the Russian for confidence, then the Russian for to doubt, then say \emph{soglashátsya} again
\end{itemize}

\begin{tcolorbox}[breakable,colback=teal!4,colframe=teal!35!black,title={Before you move on}]
What does \ru{соглашаться} mean? (\textquotedblleft{}To agree\textquotedblright{}.) Say it once more.
\end{tcolorbox}
\section[\texorpdfstring{\ru{согласие} (soglásiye)}{согласие (soglásiye)}]{\ru{согласие} (soglásiye) — agreement, consent}

\label{lesson:RU-C340-soglasie}



\begin{quote}
Before the new one: say the Russian for to doubt, then the Russian for to agree.
\end{quote}

\subsection*{You'll want to know: \ru{согласие}}
\textbf{\ru{согласие}} (\emph{soglásiye}) — \textquotedblleft{}agreement, consent\textquotedblright{}.

\begin{quote}
\ru{Родители} \ru{дали} \ru{согласие} \ru{на} \ru{свадьбу}.
\end{quote}

\begin{quote}
\emph{Rodíteli dáli soglásiye na svádbu.}
\end{quote}

\begin{quote}
The parents gave their consent to the wedding.
\end{quote}

Neuter. \emph{\ru{Согласие} \ru{на}} + accusative — consent to. \emph{\ru{В} \ru{согласии} \ru{с}} — in agreement with.

\begin{cousinweb}
\textbf{\ru{согласие}} is built from \emph{\ru{со}-} \textquotedblleft{}together\textquotedblright{} and \emph{\ru{глас}} \textquotedblleft{}voice\textquotedblright{}, which Russian took from Church Slavonic; the native form of the same Proto-Slavic word, \emph{*golsŭ}, is \emph{\ru{голос}} (EtymDB).
\end{cousinweb}

\begin{grammarlens}[title={What you've built}]
One more word for telling, explaining and describing at B1.
\end{grammarlens}

\subsection*{Your turn}
\begin{itemize}
\raggedright
  \item \emph{Say it:} \emph{soglásiye}
  \item \emph{Say it:} \emph{soglásiye}, once more
  \item \emph{Say it:} a sentence of your own with \emph{soglásiye}
  \item \emph{Recall:} say the Russian for to doubt, then the Russian for to agree, then say \emph{soglásiye} again
\end{itemize}

\begin{tcolorbox}[breakable,colback=teal!4,colframe=teal!35!black,title={Before you move on}]
What does \ru{согласие} mean? (\textquotedblleft{}Agreement, consent\textquotedblright{}.) Say it once more.
\end{tcolorbox}
\section[\texorpdfstring{\ru{несогласие} (nesoglásiye)}{несогласие (nesoglásiye)}]{\ru{несогласие} (nesoglásiye) — disagreement}

\label{lesson:RU-C340-nesoglasie}



\begin{quote}
Before the new one: say the Russian for to agree, then the Russian for agreement.
\end{quote}

\subsection*{You'll want to know: \ru{несогласие}}
\textbf{\ru{несогласие}} (\emph{nesoglásiye}) — \textquotedblleft{}disagreement\textquotedblright{}.

\begin{quote}
\ru{Он} \ru{выразил} \ru{несогласие} \ru{с} \ru{нашим} \ru{планом}.
\end{quote}

\begin{quote}
\emph{On výrazil nesoglásiye s náshim plánom.}
\end{quote}

\begin{quote}
He expressed disagreement with our plan.
\end{quote}

Neuter. \emph{\ru{Несогласие} \ru{с}} + instrumental.

\begin{cousinweb}
\textbf{\ru{несогласие}} is \emph{\ru{не}-} \textquotedblleft{}not\textquotedblright{} in front of \emph{\ru{согласие}} \textquotedblleft{}agreement\textquotedblright{}, which is built on Church Slavonic \emph{\ru{глас}} \textquotedblleft{}voice\textquotedblright{}, Proto-Slavic \emph{*golsŭ} (EtymDB).
\end{cousinweb}

\begin{grammarlens}[title={What you've built}]
One more word for telling, explaining and describing at B1.
\end{grammarlens}

\subsection*{Your turn}
\begin{itemize}
\raggedright
  \item \emph{Say it:} \emph{nesoglásiye}
  \item \emph{Say it:} \emph{nesoglásiye}, once more
  \item \emph{Say it:} a sentence of your own with \emph{nesoglásiye}
  \item \emph{Recall:} say the Russian for to agree, then the Russian for agreement, then say \emph{nesoglásiye} again
\end{itemize}

\begin{tcolorbox}[breakable,colback=teal!4,colframe=teal!35!black,title={Before you move on}]
What does \ru{несогласие} mean? (\textquotedblleft{}Disagreement\textquotedblright{}.) Say it once more.
\end{tcolorbox}
\section[\texorpdfstring{\ru{возражать} (vozrazhát)}{возражать (vozrazhát)}]{\ru{возражать} (vozrazhát) — to object}

\label{lesson:RU-C340-vozrazhat}



\begin{quote}
Before the new one: say the Russian for agreement, then the Russian for disagreement.
\end{quote}

\subsection*{You'll want to know: \ru{возражать}}
\textbf{\ru{возражать}} (\emph{vozrazhát}) — \textquotedblleft{}to object\textquotedblright{}.

\begin{quote}
\ru{Если} \ru{никто} \ru{не} \ru{возражает}, \ru{мы} \ru{начнём}.
\end{quote}

\begin{quote}
\emph{Yésli niktó ne vozrazháyet, my nachnyóm.}
\end{quote}

\begin{quote}
If nobody objects, we will begin.
\end{quote}

Imperfective; the perfective is \emph{\ru{возразить}} (\emph{\ru{возражу}, \ru{возразишь}}). \emph{\ru{Возражать} \ru{против}} + genitive. \emph{\ru{Я} \ru{не} \ru{возражаю}} — I do not mind.

\begin{cousinweb}
\textbf{\ru{возражать}} is built from \emph{\ru{воз}-} \textquotedblleft{}back, against\textquotedblright{} and the root of \emph{\ru{разить}} \textquotedblleft{}to strike\textquotedblright{} (EtymDB lists \emph{razítĭ} \textquotedblleft{}to strike\textquotedblright{}). To object is to strike back with words.
\end{cousinweb}

\begin{grammarlens}[title={What you've built}]
One more word for telling, explaining and describing at B1.
\end{grammarlens}

\subsection*{Your turn}
\begin{itemize}
\raggedright
  \item \emph{Say it:} \emph{vozrazhát}
  \item \emph{Say it:} \emph{vozrazhát}, once more
  \item \emph{Say it:} a sentence of your own with \emph{vozrazhát}
  \item \emph{Recall:} say the Russian for agreement, then the Russian for disagreement, then say \emph{vozrazhát} again
\end{itemize}

\begin{tcolorbox}[breakable,colback=teal!4,colframe=teal!35!black,title={Before you move on}]
What does \ru{возражать} mean? (\textquotedblleft{}To object\textquotedblright{}.) Say it once more.
\end{tcolorbox}
\section[\texorpdfstring{\ru{возражение} (vozrazhéniye)}{возражение (vozrazhéniye)}]{\ru{возражение} (vozrazhéniye) — an objection}

\label{lesson:RU-C340-vozrazhenie}



\begin{quote}
Before the new one: say the Russian for disagreement, then the Russian for to object.
\end{quote}

\subsection*{You'll want to know: \ru{возражение}}
\textbf{\ru{возражение}} (\emph{vozrazhéniye}) — \textquotedblleft{}an objection\textquotedblright{}.

\begin{quote}
\ru{Возражения} \ru{есть}? \ru{Нет}? \ru{Тогда} \ru{решили}.
\end{quote}

\begin{quote}
\emph{Vozrazhéniya yest? Net? Togdá reshíli.}
\end{quote}

\begin{quote}
Any objections? No? Then it is decided.
\end{quote}

Neuter. \emph{\ru{Без} \ru{возражений}} — without objection.

\begin{cousinweb}
\textbf{\ru{возражение}} is built on \emph{\ru{возразить}} \textquotedblleft{}to object\textquotedblright{} with the noun suffix \emph{-\ru{ение}}; the root is that of \emph{\ru{разить}} \textquotedblleft{}to strike\textquotedblright{} (EtymDB).
\end{cousinweb}

\begin{grammarlens}[title={What you've built}]
One more word for telling, explaining and describing at B1.
\end{grammarlens}

\subsection*{Your turn}
\begin{itemize}
\raggedright
  \item \emph{Say it:} \emph{vozrazhéniye}
  \item \emph{Say it:} \emph{vozrazhéniye}, once more
  \item \emph{Say it:} a sentence of your own with \emph{vozrazhéniye}
  \item \emph{Recall:} say the Russian for disagreement, then the Russian for to object, then say \emph{vozrazhéniye} again
\end{itemize}

\begin{tcolorbox}[breakable,colback=teal!4,colframe=teal!35!black,title={Before you move on}]
What does \ru{возражение} mean? (\textquotedblleft{}An objection\textquotedblright{}.) Say it once more.
\end{tcolorbox}
